\documentclass[10pt]{article}
\usepackage[T1]{fontenc}
\usepackage[utf8]{inputenc}
\usepackage[french]{babel}
\usepackage{amsmath,amssymb,mathrsfs,mathpazo}
\usepackage{geometry}
\geometry{a4paper,margin=2cm}
\usepackage{graphicx,array,tabularx,tabularray,colortbl,multirow,multicol,booktabs,enumitem,pifont,ulem,fancybox,diagbox,cancel}
\usepackage{pstricks,pstricks-add,pst-plot,pst-tree,pst-node,pst-text,pst-func,pst-math,pst-bspline,pst-3dplot,pst-eucl,pst-solides3d}
\let\pscalculate\undefined
\usepackage{pst-bezier}
\usepackage{tikz,pgfplots,tkz-tab}
\usetikzlibrary{arrows,arrows.meta,patterns,calc,positioning,trees,decorations.markings,shapes.geometric,shapes.misc,shapes.arrows,intersections,decorations.pathmorphing,decorations.pathreplacing,backgrounds,shadows,babel}
\pgfplotsset{compat=1.18}
\usepgfplotslibrary{fillbetween}
\usepackage[np]{numprint}
\usepackage{listings,makecell,pas-tableur,twoopt}
\let\barre\undefined
\usepackage[pstricks]{tablvar}
\definecolor{darkgreen}{RGB}{0,100,0}
\providecommand{\red}{\color{red}}
\providecommand{\blue}{\color{blue}}
\providecommand{\green}{\color{green}}
\providecommand{\black}{\color{black}}
\providecommand{\white}{\color{white}}
\providecommand{\gray}{\color{gray}}
\providecommand{\cyan}{\color{cyan}}
\providecommand{\magenta}{\color{magenta}}
\providecommand{\yellow}{\color{yellow}}
\providecommand{\violet}{\color{violet}}
\providecommand{\orange}{\color{orange}}
\providecommand{\euro}{\texteuro}
\providecommand{\EUR}{\texteuro}
\providecommand{\ieme}{\textsuperscript{e}}
\providecommand{\degres}{\ensuremath{^\circ}}
\providecommand{\degree}{\ensuremath{^\circ}}
\providecommand{\vv}[1]{\overrightarrow{#1}}
\providecommand{\Coord}[1]{#1}
\providecommand{\decosix}{}
\providecommand{\starredbullet}{\ensuremath{\bullet}}
\pagestyle{empty}
\setlength{\parindent}{0pt}
\usepackage{adjustbox}\DeclareUnicodeCharacter{2236}{:}\DeclareUnicodeCharacter{2190}{\ensuremath{\leftarrow}}
\providecommand{\R}{\mathbb{R}}
\providecommand{\N}{\mathbb{N}}
\providecommand{\Z}{\mathbb{Z}}
\providecommand{\Q}{\mathbb{Q}}
\providecommand{\D}{\mathbb{D}}
\providecommand{\C}{\mathbb{C}}
\providecommand{\e}{{\mathrm{e}}}
\providecommand{\vect}[1]{\overrightarrow{#1}}
\providecommand{\vectt}[1]{\overrightarrow{\mathrm{#1}}}
\providecommand{\vvt}[1]{\overrightarrow{\mathrm{#1}}}
\providecommand{\np}[2][]{{#2~#1}}
\providecommand{\vv}[1]{{\overrightarrow{#1}}}
\renewcommand{\d}{{\mathrm{d}}}
\providecommand{\degres}{{^\circ}}
\providecommand{\degre}{{^\circ}}
\providecommand{\blue}{\color{blue}}
\providecommand{\red}{\color{red}}
\providecommand{\enskip}{\quad}
\providecommand{\llbracket}{\lbrack\!\lbrack}
\providecommand{\rrbracket}{\rbrack\!\rbrack}
\providecommand{\euro}{{\text{€}}}
\providecommand{\no}{{\text{n°}}}
\providecommand{\interval}[2]{{[#1;#2]}}
\providecommand{\intervFO}[2]{{[#1;#2[}}
\providecommand{\dotfill}{\ldots}

\providecommand{\ds}{\displaystyle}
\providecommand{\textding}[1]{#1}
\providecommand{\Oij}{\(\mathrm{O};\vec{\imath},\vec{\jmath}\)}
\providecommand{\Oijk}{\(\mathrm{O};\vec{\imath},\vec{\jmath},\vec{k}\)}
\providecommand{\Ouv}{\(\mathrm{O};\vec{u},\vec{v}\)}
\providecommand{\intervFF}[2]{\ensuremath{[#1;#2]}}
\renewcommand{\euro}{\ifmmode\text{\texteuro}\else\texteuro\fi}
\DeclareUnicodeCharacter{00B0}{\ifmmode^{\circ}\else\textdegree\fi}
\DeclareUnicodeCharacter{20AC}{\ifmmode\text{\texteuro}\else\texteuro\fi}
% Énoncés : Éducation nationale. Transcription et corrigés : APMEP, Denis Vergès et contributeurs. Découpage et classement : CoopMaths, Sébastien Lozano, Éric Elter et contributeurs. Adaptation de la présentation : Maths Explorer.
% CC BY-SA : https://coopmaths.fr/www/about/
% https://coopmaths.fr/alea/static/bac/2022/tex/bac_2022_05_sujet2_ameriquenord_2.tex
\begin{document}
\noindent\textit{APMEP / CoopMaths. CC BY-SA. Présentation adaptée par Maths Explorer.}\par\medskip
\bigskip

\textbf{Partie A}

\medskip 

Soit $p$ la fonction définie sur l'intervalle $[-3~;~4]$ par :

\[p(x) = x^3 - 3x^2 + 5x + 1\]

\begin{enumerate}
\item Déterminer les variations de la fonction $p$ sur l'intervalle $[-3~;~4]$.
\item Justifier que l'équation $p(x) = 0$ admet dans l'intervalle $[-3~;~4]$ une unique solution qui sera notée $\alpha$.
\item Déterminer une valeur approchée du réel $\alpha$ au dixième près.
\item Donner le tableau de signes de la fonction $p$ sur l'intervalle $[-3~;~4]$. 
\end{enumerate}

\bigskip

\textbf{Partie B}

\medskip 

Soit $f$ la fonction définie sur l'intervalle $[-3~;~4]$ par :

\[f(x) = \dfrac{\text{e}^x}{1 + x^2}\]

On note $\mathcal{C}_f$ sa courbe représentative dans un repère orthogonal.

\medskip

\begin{enumerate}
\item 
	\begin{enumerate}
		\item Déterminer la dérivée de la fonction $f$ sur l'intervalle $[-3~;~4]$.
		\item Justifier que la courbe $\mathcal{C}_f$ admet une tangente horizontale au point d'abscisse 1.
	\end{enumerate}
\item  Les concepteurs d'un toboggan utilisent la courbe $\mathcal{C}_f$ comme profil d'un toboggan. Ils estiment que le toboggan assure de bonnes sensations si le profil possède au moins deux points d'inflexion.

\medskip

\begin{minipage}{0.48\linewidth}
\psset{unit=0.75cm}
\begin{pspicture*}(-4,-1.95)(4.25,4.1)
\psgrid[gridlabels=0pt,subgriddiv=1,gridwidth=0.15pt](-4,0)(4,4)
\psaxes[linewidth=1.25pt,labelFontSize=\scriptstyle](0,0)(-4,-0)(4.25,4.1)
\psplot[plotpoints=2000,linewidth=1.25pt,linecolor=blue]{-3}{4}{2.71828 x exp x dup mul 1 add div}
\rput(0,-1){Représentation de la courbe $\mathcal{C}_f$}
\uput[ul](3,2){\blue \small $\mathcal{C}_f$}
\end{pspicture*}
\end{minipage} \hfill
\begin{minipage}{0.48\linewidth}
\psset{unit=0.75cm}
\begin{pspicture}(-4,-1.95)(4.25,4.1)

\def\tobo{\psplot[plotpoints=2000,linewidth=1.25pt]{-3}{4}{2.71828 x exp x dup mul 1 add div}}
\psplot[plotpoints=2000,linewidth=1.25pt]{-3}{4}{2.71828 x exp x dup mul 1 add div}
\rput(-0.6,0.6){\tobo}
\psline(-3,0)(-3.6,0.6)\psline(4,3.2)(3.4,3.8)
\pscustom[fillstyle=solid,fillcolor=lightgray]
{\psplot[plotpoints=2000,linewidth=1.25pt,linecolor=red]{-3}{4}{2.71828 x exp x dup mul 1 add div}
\psline(4,3.2)(4,0)(-2,0)}
\rput(0,-1){Vue de profil du toboggan}
\end{pspicture}
\end{minipage}

	\begin{enumerate}
		\item D'après le graphique ci-dessus, le toboggan semble-t-il assurer de bonnes sensations ? Argumenter.
		\item On admet que la fonction $f''$, dérivée seconde de la fonction $f$, a pour expression pour tout réel $x$ de l'intervalle $[-3~;~4]$ :

		\[f''(x) = \dfrac{p(x)(x - 1)\text{e}^x}{\left(1 + x^2\right)^3}\]

où $p$ est la fonction définie dans la partie A.

En utilisant l'expression précédente de $f''$, répondre à la question : \og le toboggan assure-t-il de bonnes sensations ? \fg. Justifier.
	\end{enumerate}
\end{enumerate}

\bigskip

\end{document}
